\section{Extended Related Work}
\label{sec:app:related_works}
\paragraph{Detecting backdoors.}
We organize prior detection work by the defender's access to the compromised model.
\textit{Black-box} methods audit a deployed model through its outputs, typically via a trusted auditor LLM \citep{petri2025, petri2026v2}; some compromised models will even confess to being modified under the right elicitation \citep{sheshadri_auditbench}.
\textit{White-box} methods exploit memorization of training-time data: \citet{bullwinkel_trigger_2026} extract poisoned samples directly from the weights, but their threat model assumes a single fixed trigger phrase prepended to all prompts and does not extend to dispersed or variable-position triggers \citep{hubinger_sleeper_2024, huang-etal-2024-composite}.
\textit{Gray-box} methods, which use internal representations without requiring the training corpus, remain comparatively underexplored. \citet{zeng_beear_2024} observe that backdoor triggers induce a relatively uniform drift in the embedding space and use this for mitigation rather than detection.
